\section*{Capítulo \ref{ch:b2:anthropogenic-climate-change} --- A mudança climática antrópica}

\begin{solution}{exo:b2:anthropogenic-climate-change:1}
(a) $3.75 - 0.60 + 0.22 + 0.45 + 0.11 = \qty{3.93}{W.m^{-2}}$, os ajustes somando
\qty{0.18}{W.m^{-2}}. (b) $3.93/3.71 = 1.059$: \qty{5.9}{\percent} a mais. (c) O ajuste de temperatura
troposférico, $-0.60$: uma troposfera mais quente emite mais infravermelho para o espaço, o que retira
parte da energia adicional antes que a superfície tenha respondido.
\end{solution}

\begin{solution}{exo:b2:anthropogenic-climate-change:2}
(a) $4\pi(\num{6.371e6})^2 = \qty{5.10e14}{m^2}$. (b) Em 47 anos, $\num{434.9e21}/(47 \times
\num{3.156e7} \times \num{5.10e14}) = \qty{0.575}{W.m^{-2}}$; em 12 anos, $\num{152.4e21}/(12
\times \num{3.156e7} \times \num{5.10e14}) = \qty{0.789}{W.m^{-2}}$, os 0.57 e 0.79 da
avaliação. (c) $0.91 \times 434.9 = \qty{396}{ZJ}$.
\end{solution}

\begin{solution}{exo:b2:anthropogenic-climate-change:3}
(a) $-3.22 + 1.30 + 0.35 + 0.42 - 0.01 = \num{-1.16}$, a líquida. (b) $3.93/1.16 =
\qty{3.39}{\celsius}$; $3.93/3.22 = \qty{1.22}{\celsius}$. (c) $3.39/1.22 = 2.78$: as retroalimentações
quase triplicam o aquecimento de Planck.
\end{solution}

\begin{solution}{exo:b2:anthropogenic-climate-change:4}
(a) Um petagrama de carbono equivale a $44.009/12.011 = 3.664$~Gt de dióxido de carbono: $0.45 \times 3.664 =
\qty{1.65}{\celsius}$ por \qty{1000}{PgC}, e o intervalo torna-se 0.989 a 2.31. (b) $2390 \times
12.011/44.009 = \qty{652}{PgC}$. (c) $0.45 \times 2.39 = \qty{1.08}{\celsius}$.
\end{solution}

\begin{solution}{exo:b2:anthropogenic-climate-change:5}
(a) $g = 159/398 = 0.400$; $T = (398/\num{5.670e-8})^{1/4} = \qty{289}{K}$. (b) $-0.6 \times 4
\times \num{5.670e-8} \times 288^3 = \qty{-3.25}{W.m^{-2}.K^{-1}}$; a \qty{298}{K}, $-3.60$.
(c) Menos sensível: a resposta cresce como $T^3$, de modo que um planeta mais quente se livra de cada grau adicional mais
facilmente, sendo iguais as outras retroalimentações.
\end{solution}

\begin{solution}{exo:b2:anthropogenic-climate-change:6}
(a) $1.85 - 0.50 = \qty{1.35}{W.m^{-2}.\celsius^{-1}}$. (b) $\sqrt{0.32^2 + 0.20^2} = 0.38$.
(c) A soma é muito mais bem conhecida que qualquer dos termos: seus erros são fortemente anticorrelacionados. Onde
a alta troposfera aquece mais, ela também contém mais vapor, de modo que uma \omterm{def:b2:anthropogenic-climate-change:feedbacks}{retroalimentação do gradiente vertical}
mais negativa vem com uma \omterm{def:b2:anthropogenic-climate-change:feedbacks}{retroalimentação do vapor d'água} mais positiva.
\end{solution}

\begin{solution}{exo:b2:anthropogenic-climate-change:7}
(a) $\lambda = (0.59 - 2.20)/1.03 = \qty{-1.56}{W.m^{-2}.\celsius^{-1}}$; $3.93/1.56 =
\qty{2.52}{\celsius}$. (b) $3.93/(1.56 - 0.5) = 3.93/1.06 = \qty{3.71}{\celsius}$. (c) $3.93/1.46 =
\qty{2.69}{\celsius}$. A avaliação, com distribuições completas, encontra 3.5 e \qty{2.7}{\celsius}.
\end{solution}

\begin{solution}{exo:b2:anthropogenic-climate-change:8}
(a) $0.617/3682 = \qty{1.68e-4}{\celsius^{-1}}$. (b) $0.113 \times 0.396 = \qty{0.0447}{m}$, ou
\qty{44.7}{mm}, perto dos \qty{47.5}{mm} do balanço. (c) $\num{2e-4} \times 0.2 \times 700 =
\qty{0.028}{m}$, \qty{28}{mm}.
\end{solution}

\begin{solution}{exo:b2:anthropogenic-climate-change:9}
(a) $\sum xy = 0.311$, $\sum x^2 = 0.300$: $\beta = 1.04$. (b) O intervalo de \qty{95}{\percent} é
$1.04 \pm 1.96 \times 0.08$, 0.88 a 1.20: ele exclui 0, logo a resposta é detectada, e contém
1, logo a mudança observada é coerente com a simulada. (c) O intervalo 0.34 a 0.66 ainda
exclui 0 (detectada), mas exclui 1: o modelo superestimaria a resposta por um fator de cerca
de dois.
\end{solution}

\begin{solution}{exo:b2:anthropogenic-climate-change:10}
(a) $\sum g^2 = 1.20$, $\sum a^2 = 0.18$, $\sum ga = -0.46$, $\sum gy = 0.74$, $\sum ay = -0.28$.
As equações normais são $1.20\,\beta_g - 0.46\,\beta_a = 0.74$ e $-0.46\,\beta_g + 0.18\,\beta_a
= -0.28$; seu determinante é $1.20 \times 0.18 - 0.46^2 = 0.0044$, e $\beta_g = (0.74 \times
0.18 - 0.46 \times 0.28)/0.0044 = 1.00$, $\beta_a = (-1.20 \times 0.28 + 0.46 \times 0.74)/0.0044 =
1.00$. (b) Gases de efeito estufa $+\qty{0.80}{\celsius}$, aerossóis $-\qty{0.30}{\celsius}$. (c) O
determinante é pequeno porque $g$ e $a$ são quase proporcionais (sua correlação é $-0.99$): um
pouco de ruído em $y$ pode ser absorvido elevando os dois fatores de escala juntos, de modo que cada um é mal
determinado, embora seu efeito combinado seja bem conhecido.
\end{solution}

\begin{solution}{exo:b2:anthropogenic-climate-change:11}
(a) $(0.93 - 0.1)/0.45 \times 1000 = \qty{1844}{Gt}$; menos 210 e $26 \times 0.93 = 24$:
\qty{1610}{Gt}. (b) $0.93 - x = 0.45 \times (1350 + 210 + 24)/1000 = 0.713$, logo $x =
\qty{0.22}{\celsius}$: aproximadamente o limite superior dos 0.1 a \qty{0.2}{\celsius} avaliados. (c) O
orçamento depende do aquecimento não devido ao CO$_2$, $x$, apenas por meio de $-x/\text{TCRE}$: cada
\qty{0.1}{\celsius} custa $0.1/0.45 \times 1000 = \qty{222}{Gt}$, qualquer que seja o limite.
\end{solution}

\begin{solution}{exo:b2:anthropogenic-climate-change:12}
(a) $3.93/1.81 = \qty{2.17}{\celsius}$, $3.93/0.51 = \qty{7.71}{\celsius}$, $3.93/1.16 =
\qty{3.39}{\celsius}$. (b) O lado baixo tem $3.39 - 2.17 = \qty{1.22}{\celsius}$ de comprimento, o lado alto
$7.71 - 3.39 = \qty{4.32}{\celsius}$: a função $1/|\lambda|$ é convexa, de modo que passos iguais de
$\lambda$ em direção a zero elevam a ECS cada vez mais. (c) Um $\lambda$ fracamente negativo dá uma ECS muito
grande, e pequenas variações de $\lambda$ perto de zero são difíceis de excluir; uma ECS baixa exige um $\lambda$
fortemente negativo, que a \omterm{def:b2:anthropogenic-climate-change:feedbacks}{resposta de Planck} limita e que o aquecimento observado contradiz.
\end{solution}

\begin{solution}{pb:b2:anthropogenic-climate-change:1}
\textbf{1.} $2.16 + 0.54 + 0.21 + 0.41 = \qty{3.32}{W.m^{-2}}$.
\textbf{2.} $-0.22 - 0.84 = \qty{-1.06}{W.m^{-2}}$.
\textbf{3.} $3.32 + 0.47 + 0.05 - 1.06 - 0.20 + 0.08 + 0.06 = \qty{2.72}{W.m^{-2}}$, o total
avaliado.
\textbf{4.} $3.93\log_2(410/278) = \qty{2.20}{W.m^{-2}}$, a menos de \qty{2}{\percent} de 2.16.
\textbf{5.} $3.93\log_2(427.35/278) = \qty{2.44}{W.m^{-2}}$.
\textbf{6.} O termo aerossóis--nuvens, $-1.45$ a $-0.25$, com \qty{1.20}{W.m^{-2}} de largura, contra 1.52 para
o total: ele carrega a maior parte da incerteza do total porque é grande e mal
determinado.
\textbf{7.} $0.79 \times \num{5.10e14} \times 12 \times \num{3.156e7} = \qty{1.53e23}{J}$,
\qty{153}{ZJ}.
\textbf{8.} $0.91 \times 153 = \qty{139}{ZJ}$.
\textbf{9.} $\lambda = (0.59 - 2.20)/1.03 = \qty{-1.56}{W.m^{-2}.\celsius^{-1}}$.
\textbf{10.} $3.93/1.56 = \qty{2.52}{\celsius}$.
\textbf{11.} $3.93 \times 1.03/2.20 = \qty{1.84}{\celsius}$.
\textbf{12.} $3.93/1.16 = \qty{3.39}{\celsius}$, mais alta em \qty{0.87}{\celsius}.
\textbf{13.} O aquecimento até agora foi mais forte onde as retroalimentações são mais negativas; à medida que o
padrão evolui, o parâmetro de retroalimentação torna-se menos negativo, de modo que a sensibilidade de longo prazo
excede a efetiva.
\textbf{14.} $0.45 \times 44.009/12.011 = \qty{1.65}{\celsius}$ por \qty{1000}{PgC}.
\textbf{15.} 0.989 a \qty{2.31}{\celsius} por \qty{1000}{PgC}.
\textbf{16.} $2390 \times 12.011/44.009 = \qty{652}{PgC}$.
\textbf{17.} $0.45 \times 2.39 = \qty{1.08}{\celsius}$.
\textbf{18.} Dentro do arredondamento, o dióxido de carbono sozinho explica todo o aquecimento humano: o
aquecimento pelos outros gases de efeito estufa e o resfriamento pelos aerossóis quase se compensam.
\textbf{19.} $1.5 - 1.07 = \qty{0.43}{\celsius}$.
\textbf{20.} $(0.43 - 0.1)/0.45 \times 1000 = \qty{733}{Gt}$ de dióxido de carbono.
\textbf{21.} $733 - 210 = \qty{523}{Gt}$.
\textbf{22.} $26 \times 0.43 = \qty{11.2}{Gt}$; $523 - 11.2 = \qty{512}{Gt}$.
\textbf{23.} Ele cai $0.1/0.45 \times 1000 = \qty{222}{Gt}$, para \qty{290}{Gt}: os
\qty{220}{Gt} pelos quais a avaliação diz que as escolhas sobre os outros fatores deslocam o orçamento.
\textbf{24.} Calculado a partir da TCRE, o \omterm{def:b2:anthropogenic-climate-change:budget}{orçamento de carbono restante} para \qty{1.5}{\celsius} a partir do
início de 2020 é de \textbf{\qty[text-series-to-math=true]{512}{Gt}} de dióxido de carbono, a menos de
\qty{3}{\percent} dos \qty{500}{Gt} avaliados com uma probabilidade de \qty{50}{\percent}.
\end{solution}
